\section{Hoare Logic for IMP}
\label{sec:imp}

We build a Hoare logic for a small while-language IMP, calibrated for
the conceptual point of this paper: showing how the forcing-style
propositional layer lifts to a sound Hoare logic, and proving the
\texttt{while} rule both by meta-level induction on the step index
and by an explicit appeal to the internal Löb rule.  Exhibiting both
proofs side by side makes precise the claim that ``Löb's rule
\emph{is} induction on the step index.''  IMP has no higher-order
features, no allocation, and no concurrency, so the step indices are
not yet doing technical work.

IMP commands and their small-step semantics are standard:
\begin{lstlisting}
Inductive cmd : Type :=
  | CSkip | CAssign : var -> aexp -> cmd
  | CSeq : cmd -> cmd -> cmd
  | CIf : bexp -> cmd -> cmd -> cmd
  | CWhile : bexp -> cmd -> cmd.
\end{lstlisting}
\noindent
Each small step corresponds to one decrement of the step index
tracked by the wp predicate.

\paragraph{Step-indexed wp and the Hoare triple.}
The wp predicate is the step-indexed analogue of the Dijkstra wp:
$\mathtt{wp\_at}\, n\, c\, Q\, s$ says that running $c$ from $s$ for at
most $n$ steps is safe and any final state satisfies $Q$.  In IMP,
every non-\texttt{CSkip} command can step, so the safety condition is
automatic and we do not record it explicitly.
\begin{lstlisting}
Fixpoint wp_at (n : nat) (c : cmd)
               (Q : state -> Prop) (s : state) : Prop :=
  match n with
  | O => True
  | S k => (c = CSkip -> Q s) /\
           (forall c' s', step (c, s) (c', s') -> wp_at k c' Q s')
  end.
\end{lstlisting}
\noindent
The predicate is downward-closed in $n$ and contravariantly monotone
in $Q$, and lifts to an \texttt{iProp} via \texttt{MkProp}.  A Hoare
triple is the \texttt{iProp} saying ``for every $s$ with $P\, s$, the
wp of $c$ targeting $Q$ holds.''  Validity, written
$\mathtt{hoare\_valid}\, P\, c\, Q$, is the entailment
$\top \vdash \mathtt{hoare}\, P\, c\, Q$. Equivalently, by
\texttt{hoare\_valid\_alt}, it is the meta-level statement ``for every
$n$ and $s$ with $P\, s$, the $n$-step-bounded wp holds.''

Partial-correctness Hoare logic has five structural rules: skip,
assignment, sequencing, conditional, and consequence.  Each becomes a
small lemma about \texttt{wp\_at}, provable by routine unfolding at
$n+1$ and inversion on the available step.  Sequencing requires a
small auxiliary lemma \texttt{wp\_seq} decomposing the wp of
\texttt{CSeq c1 c2} into the wp of \texttt{c1} targeting the wp of
\texttt{c2}, with one use of \texttt{wp\_at\_mono} to align indices.
None of these five rules uses the propositional layer, and they
could be stated entirely at the meta level.  The propositional layer
becomes useful only for \texttt{while}.

\paragraph{The \texttt{while} rule, meta-level induction.}
Soundness of
$\{P \wedge b\}\, c\, \{P\} \;\vdash\; \{P\}\, \mathtt{while}\, b\, c\, \{P \wedge \neg b\}$
says that a sound loop body gives a sound loop, and the proof goes by
induction on the step index.  The base case is trivial. The inductive
step unfolds the wp at $n+1$, case-analyses on the step, and uses the
induction hypothesis to discharge the recursive \texttt{while}
obligation:
\begin{lstlisting}
Lemma hoare_while P b c :
  hoare_valid (fun s => P s /\ beval s b = true) c P ->
  hoare_valid P (CWhile b c)
              (fun s => P s /\ beval s b = false).
\end{lstlisting}

\paragraph{The same rule, via the internal Löb.}
The same theorem also follows from the internal Löb rule of
Section~\ref{sec:prop} without any explicit induction on $n$.  We
factor the goal $\top \vdash \mathtt{hoare}\, P\, (\mathtt{CWhile}\,
b\, c)\, Q$ through the iProp $\triangleright \varphi \to \varphi$ for
$\varphi := \mathtt{hoare}\, P\, (\mathtt{CWhile}\, b\, c)\, Q$. We
prove the intermediate entailment directly, by a finite, non-recursive
argument that case-analyses on the depth supplied to the
\texttt{iImpl}, and close with \texttt{apply iLob}:
\begin{lstlisting}
Lemma hoare_while_iLob P b c :
  hoare_valid (fun s => P s /\ beval s b = true) c P ->
  hoare_valid P (CWhile b c) (fun s => P s /\ beval s b = false).
\end{lstlisting}
\noindent
The intermediate entailment contains no induction on $n$.  All recursive content
is in \texttt{iLob}'s proof, written once in
Section~\ref{sec:prop} rather than at every soundness proof.  The six
rules (skip, assign, seq, if, while, consequence) are packaged as a
single theorem \texttt{imp\_hoare\_rules}.  In the forcing-style
semantics, soundness of Hoare logic over IMP is structurally
identical to the soundness statement using the standard direct
semantics: at the first-order level the forcing reading is a
different \emph{presentation} of the same theorems.
Section~\ref{sec:lam} moves to a setting where the re-presentation
produces a different proof obligation.
